%% parts/p04/machines-effective-procedures.tex --- Phase-4 migration of
%% v1.1 Part V, ch:turing-machines, source lines 7497--7849.
%% Body = definitions, statements, warnings, and method descriptions only;
%% every source proof/sketch is represented in the parallel annex fragment.
%% Depth additions: D-P5-001 and D-P5-002.  Apparatus: A-045.
%% Source lines 7494--7496 are the part heading already owned by the P4 driver;
%% they are accounted for there and intentionally not duplicated in this file.

\chapter{Machines and Effective Procedures}
\label{sec:p04-machines-effective-procedures}

Deterministic single-tape Turing machines are the primary computation model
in this part.  Alternative standard formalisms enter only after this model and
its numerical conventions have been fixed.  Historical sources are Church
and Turing \cite{Church1936,Turing1937}; systematic references include
\cite{Shoenfield1967,BoolosBurgessJeffrey2007}.

\section{Deterministic single-tape Turing machines}
\label{sec:p04-machines}

\begin{hbdefinition}{def:p04-turing-machine}{Turing machine}
A \emph{deterministic single-tape Turing machine} (TM) is a tuple
\[
  \mathcal T=(Q,\mathcal A_{\mathrm{in}},\mathcal A_{\mathrm{tape}},
  \delta,q_0,q_{\mathrm{acc}},q_{\mathrm{rej}})
\]
such that:
\begin{enumerate}
  \item \(Q\) is a finite set of states;
  \item \(\mathcal A_{\mathrm{in}}\) is a finite input alphabet not
        containing the blank symbol \(\sqcup\);
  \item \(\mathcal A_{\mathrm{tape}}\) is a finite tape alphabet with
        \(\mathcal A_{\mathrm{in}}\cup\{\sqcup\}
        \subseteq\mathcal A_{\mathrm{tape}}\);
  \item \(q_0,q_{\mathrm{acc}},q_{\mathrm{rej}}\in Q\), with
        \(q_{\mathrm{acc}}\ne q_{\mathrm{rej}}\); and
  \item
  \[
    \delta:(Q\setminus\{q_{\mathrm{acc}},q_{\mathrm{rej}}\})
    \times\mathcal A_{\mathrm{tape}}
    \longrightarrow Q\times\mathcal A_{\mathrm{tape}}\times\{L,R\}.
  \]
\end{enumerate}
The alphabet symbols in this tuple are machine components, not formula
contexts or arithmetical-hierarchy classes.  The tape is infinite in both
directions, has finite nonblank support at each finite time, and has one head.
\index{Turing machine!deterministic single-tape}
\end{hbdefinition}

\begin{hbdefinition}{def:p04-configuration-computation}{Configuration and computation}
A \emph{configuration} of \(\mathcal T\) is \(C=(q,T,h)\), where
\(q\in Q\), \(h\in\mathbb Z\), and
\(T:\mathbb Z\to\mathcal A_{\mathrm{tape}}\) has finite nonblank support.  If
\(q\notin\{q_{\mathrm{acc}},q_{\mathrm{rej}}\}\) and
\(\delta(q,T(h))=(q',a,d)\), let \(T'\) agree with \(T\) away from
\(h\), put \(T'(h)=a\), and set \(h'=h-1\) for \(d=L\) and
\(h'=h+1\) for \(d=R\).  Then, and only then,
\[
  (q,T,h)\longrightarrow_{\mathcal T}(q',T',h').
\]
A halting-state configuration has no successor; write
\(\longrightarrow_{\mathcal T}^{*}\) for reflexive transitive closure.

For \(w\in\mathcal A_{\mathrm{in}}^*\) of length \(n\), let \(T_w(i)\)
be the \(i\)-th symbol of \(w\) for \(0\le i<n\), and \(\sqcup\)
otherwise.  The initial configuration is \((q_0,T_w,0)\).  A
\emph{computation on \(w\)} is the unique maximal finite or infinite
sequence from this configuration following \(\longrightarrow_{\mathcal T}\).
It halts exactly on reaching a halting state and accepts or rejects according
to which one it reaches.  Thus the empty input starts with the head scanning
a blank in cell \(0\).
\index{Turing machine!configuration}
\index{Turing machine!computation}
\end{hbdefinition}

\begin{hbconvention}{def:p04-numerical-coding}{Binary numerical and tuple coding}
For \(n\in\N\), \(\operatorname{bin}(0)\) is the one-symbol word \(0\),
and for \(n>0\), \(\operatorname{bin}(n)\) is the ordinary base-two
expansion without leading zeros, most significant bit first.  Use
\[
  \operatorname{pair}_{\N}(a,b)
  =\frac{(a+b)(a+b+1)}2+b,
\]
a numerical pairing map distinct in notation and type from the set-theoretic
ordered pair \(\langle a,b\rangle\).  For each fixed arity set
\[
 \operatorname{code}_0()\defeq0,
 \qquad \operatorname{code}_1(n)\defeq n,
\]
and, for \(k\geq1\),
\[
 \operatorname{code}_{k+1}(n_0,\ldots,n_k)
 \defeq
 \operatorname{pair}_{\N}\bigl(n_0,
   \operatorname{code}_k(n_1,\ldots,n_k)\bigr).
\]
The type \(\N^k\), rather than the code, supplies the arity.  The numerical
input for \(\vec n\in\N^k\) is
\(\operatorname{bin}(\operatorname{code}_k(\vec n))\), placed from cell
\(0\) with its most significant bit first.
\index{binary numerical coding}
\index{tuple coding}
\end{hbconvention}

\begin{hbconvention}{def:p04-numerical-output}{Numerical output}
A machine used to compute a numerical function has
\(\{0,1\}\subseteq\mathcal A_{\mathrm{in}}\).  For a halting tape \(T\), put
\[
  m\defeq\min\{i\in\N:T(i)=\sqcup\}.
\]
Finite nonblank support makes \(m\) defined.  The prefix of length \(m\)
is a \emph{valid numerical output} exactly when it is
\(\operatorname{bin}(y)\) for a unique \(y\in\N\).  Hence the empty word,
a word containing a nondigit, and a nonzero binary word with a leading zero
are invalid.  Cells outside the prefix and the final head position do not
affect the output.  A numerical-function machine must halt in
\(q_{\mathrm{acc}}\) with a valid numerical output.
\index{Turing machine!numerical output}
\end{hbconvention}

\begin{hbdefinition}{def:p04-partial-computable}{Partial computable function}
A partial function \(f:\N^k\rightharpoonup\N\) is \emph{partial
computable} iff some Turing machine satisfies both of the following conditions
for every \(\vec n\in\N^k\): when \(f(\vec n)\) is defined, the computation
on \(\operatorname{bin}(\operatorname{code}_k(\vec n))\) reaches
\(q_{\mathrm{acc}}\) with output
\(\operatorname{bin}(f(\vec n))\); when \(f(\vec n)\) is undefined, that
computation is infinite.  A total partial computable function is
\emph{computable}.  The arrow \(\rightharpoonup\) permits undefined values;
divergence is not a third numerical output.
\index{partial computable function}
\index{computable function}
\end{hbdefinition}

Definition~\ref{def:p04-partial-computable} is the repaired R-035 definition:
it contains one definiendum--definiens biconditional and then two cases, not a
nested or doubled ``if and only if.''

\begin{hbmethod}{def:p04-coding-procedures}{Effective encoding and decoding}
The fixed codes introduce no hidden computation model.  To invert
\(\operatorname{pair}_{\N}\) at \(z\), boundedly search for the greatest
\(s\le z\) with \(s(s+1)/2\le z\), then put
\(b=z-s(s+1)/2\) and \(a=s-b\).  Given the arity, iteration decodes a tuple.
Repeated division by two produces a binary word, while the recurrence
\(r\mapsto2r+b\) decodes a word from left to right.  Each is a terminating
Turing-machine procedure.
\end{hbmethod}

\begin{hbexample}{def:p04-unary-successor}{Unary successor}
Because \(\operatorname{code}_1(n)=n\), a successor machine scans the input
\(\operatorname{bin}(n)\) rightward, propagates a carry left through trailing
\(1\)s, and writes \(1\) at the first \(0\).  If the carry passes the most
significant bit, it shifts the resulting zero block right, writes \(1\) in
cell \(0\), and leaves the following cell blank, so \(111\) becomes
\(1000\).  The invariant says that the
processed suffix is the correct low-order suffix of \(n+1\), while the state
records whether a carry remains.  The final prefix is an output word, not a
second tuple code.
\end{hbexample}

\begin{hbwarning}{def:p04-formalism-complexity}{Machine formalism versus complexity measure}
Changing from one tape to several does not change which partial functions are
computable, but it can change running time by a polynomial factor.  A
computability claim and a complexity claim therefore require different
invariance arguments; see \cite{BoolosBurgessJeffrey2007} for the multi-tape
argument.
\end{hbwarning}

\section{Primitive recursion}
\label{sec:p04-primitive-recursion}

\begin{hbdefinition}{def:p04-primitive-recursive-function}{Primitive recursive function}
The \emph{primitive recursive functions} form the least class of total
functions on finite powers of \(\N\) containing, for every \(k\in\N\), the
zero function \(Z_k:\N^k\to\N\) (including \(Z_0()=0\)); successor
\(S(n)=n+1\); and every projection
\(P_i^k(n_1,\ldots,n_k)=n_i\) for \(k\ge1\) and \(1\le i\le k\).  The class
is closed under composition
\[
  f(\vec x)=g(h_1(\vec x),\ldots,h_m(\vec x))
\]
and, for every \(k\in\N\), under primitive recursion from
\(g:\N^k\to\N\) and \(h:\N^{k+2}\to\N\):
\[
  f(\vec x,0)=g(\vec x),
  \qquad
  f(\vec x,n+1)=h(\vec x,n,f(\vec x,n)).
\]
\index{primitive recursive function}
\end{hbdefinition}

\begin{hbdefinition}{def:p04-primitive-recursive-relation}{Primitive recursive relation}
A relation \(R\subseteq\N^k\) is \emph{primitive recursive} iff its
characteristic function \(\operatorname{char}_R:\N^k\to\{0,1\}\), equal to \(1\) exactly
on tuples in \(R\), is primitive recursive.  This term classifies a relation,
not a representing arithmetic formula.  The distinction is required because
the later coding predicates for formulas and derivations are relations, while
Definition~\ref{def:p04-primitive-recursive-function} concerns functions.
\index{primitive recursive relation}
\end{hbdefinition}

\begin{hbtheorem}{thm:p04-pr-functions-computable}{sketched=annex:p04-machines.1}{Primitive recursive functions are computable}
Every primitive recursive function is computable and total.
\end{hbtheorem}

\begin{hbexample}{def:p04-addition-multiplication-pr}{Addition and multiplication}
List addition before multiplication by the following recursive presentations:
\[
\begin{aligned}
  x+0&=x, & x+(y+1)&=S(x+y),\\
  x\times0&=0, & x\times(y+1)&=(x\times y)+x.
\end{aligned}
\]
In the second line, \(+\) denotes the addition function presented in the first
line.
The Ackermann function is total and computable but not primitive recursive;
its domination proof is outside this entry and may be found in
\cite{BoolosBurgessJeffrey2007}.
\end{hbexample}

\section[Minimization and Formalisms]{Unbounded minimization and equivalent formalisms}
\label{sec:p04-partial-recursion}

\begin{hbconvention}{def:p04-kleene-equality}{Kleene equality}
For expressions denoting possibly undefined values, \(a\simeq b\) means that
both are undefined or both are defined and equal.  For partial functions
\(f,g:\N\rightharpoonup\N\), \(f\simeq g\) abbreviates
\(\forall x\in\N\,f(x)\simeq g(x)\).  This convention is placed here, ahead
of its position in the migrated source, because the new normal-form entry is
its first handbook use.
\index{Kleene equality}
\end{hbconvention}

\begin{hbdefinition}{def:p04-partial-recursive}{Unbounded minimization and partial recursive functions}
For a partial map \(g:\N^{k+1}\rightharpoonup\N\), its \emph{unbounded
minimization} is the partial map
\[
  (\mu g)(\vec x)=\mu y\,[g(\vec x,y)=0].
\]
This value is defined as the least \(y\) for which \(g(\vec x,y)\) is defined
and equals \(0\), provided \(g(\vec x,z)\) is defined and nonzero for every
\(z<y\); it is undefined if an earlier required computation diverges or no
such zero is reached.

For each arity, the \emph{partial recursive functions}, also called the
\emph{\(\mu\)-recursive functions}, are the least class of
partial maps between finite powers of \(\N\) that contains the zero,
successor, and projection functions and is closed under composition, primitive
recursion, and unbounded minimization.  Composition is defined only when all
required inner values and then the outer value are defined; a recursively
defined value is defined only when its base value and every preceding value
needed by the step clause are defined.  This is a class of partial numerical
functions, not a class of formulas or machine descriptions.  It strictly adds
the possibility of divergence to the total class in
Definition~\ref{def:p04-primitive-recursive-function}; that distinction is
needed by the normal-form and machine-equivalence theorems below.
\index{unbounded minimization}
\index{partial recursive function}
\index{mu-recursive function@\(\mu\)-recursive function}
\end{hbdefinition}

\begin{hbtheorem}{thm:p04-kleene-normal-form}{cited=Kleene1952}{Kleene normal form}
For each \(k\ge1\), there are a primitive recursive relation
\(\mathsf T_k(e,x_1,\ldots,x_k,y)\) and a primitive recursive function
\(\mathsf U(y)\) such that every partial recursive
\(f:\N^k\rightharpoonup\N\) has an index \(e\), in the fixed computation
coding of this normal-form theorem, satisfying
\[
  f(x_1,\ldots,x_k)
  \simeq
  \mathsf U\bigl(\mu y\,
    \mathsf T_k(e,x_1,\ldots,x_k,y)\bigr),
\]
where the minimization selects the least \(y\) satisfying \(\mathsf T_k\).
Conversely, every partial function displayed in this form is partial
recursive.  Here \(\mathsf T_k\) is the relation saying that \(y\) is a valid
finite accepting-computation code under the computation coding fixed for this
normal form, and \(\mathsf U\) extracts its numerical output; it is not the
bounded object-language computation formula introduced later for
arithmetization.
\end{hbtheorem}

\begin{hbtheorem}{thm:p04-tm-partial-recursive}{cited=Kleene1952}{Turing-computable iff partial recursive}
For every \(k\ge1\) and every partial map
\(f:\N^k\rightharpoonup\N\), the following are equivalent:
\begin{enumerate}
  \item \(f\) is computed by a Turing machine under
        Definitions~\ref{def:p04-turing-machine} and
        \ref{def:p04-partial-computable};
  \item \(f\) is partial recursive in the closure sense of
        Definition~\ref{def:p04-partial-recursive}.
\end{enumerate}
\end{hbtheorem}

The equivalence is at the level of partial numerical functions, after the
input/output encodings are fixed.  Historical equivalences among general recursion,
Turing computation, and lambda definability are in
\cite{Kleene1936,Turing1937Lambda}; \cite{Kleene1952} gives the modern
partial-function development.

The index \(e\) in Theorem~\ref{thm:p04-kleene-normal-form} belongs to that
cited normal-form coding; it is not yet being identified with an index in the
handbook's bespoke program numbering.  The next chapter fixes the single
effective program numbering used by every index-dependent construction in this
part and verifies its enumeration and parameter-specialization properties.  No
second handbook program numbering or unrelated compiler is introduced.

\begin{hbdefinition}{def:p04-shepherdson-sturgis-urm}{Shepherdson--Sturgis Unlimited Register Machine}
The \emph{Shepherdson--Sturgis Unlimited Register Machine} (URM) has registers
\(1,2,\ldots\), each storing a natural number, and a finite program mentioning
only finitely many registers.  Its six instruction forms are
\(P(n)\), increment register \(n\); \(D(n)\), decrement a nonzero register
\(n\); \(O(n)\), set register \(n\) to zero; \(C(m,n)\), copy register \(m\)
to register \(n\), with \(m\ne n\); \(J[E_1]\), transfer unconditionally;
and \(J(m)[E_1]\), take exit \(E_1\) exactly when register \(m\) is zero and
otherwise take the normal next-line exit.  Execution begins at line \(1\) and
halts on reaching the nonexistent line after the last instruction.

To compute \(f:\N^n\rightharpoonup\N\), choose distinct register addresses
\(x_1,\ldots,x_n,y\) and a bound \(N\) containing them.  Initialize the
\(x_i\)-registers with the inputs.  If \(f(\vec x)\) is undefined, the program
does not stop; if it is defined, the program stops with register \(y\)
containing \(f(\vec x)\), while registers \(1,\ldots,N\) other than \(y\)
retain their initial contents.  Registers above \(N\) may be scratch space.
This name denotes that exact formalism and function convention, not an
unspecified generic register machine; the distinction is required for its row in
Table~\ref{tab:p04-machine-equivalences} \cite{ShepherdsonSturgis1963}.
\index{Unlimited Register Machine!Shepherdson--Sturgis}
\index{register machine!Shepherdson--Sturgis}
\end{hbdefinition}

\begin{landscape}
\begingroup\small\setlength{\tabcolsep}{3pt}
\begin{longtable}{@{}p{1.10in}p{2.25in}p{1.85in}p{1.84in}@{}}
\caption{Machine models, recursive-function formalisms, and their exact
evidence relationships (A-045).}
\label{tab:p04-machine-equivalences}\\
\toprule
Model & Object and function convention & Verified relationship & Evidence \\
\midrule
\endfirsthead
\toprule
Model & Object and function convention & Verified relationship & Evidence \\
\midrule
\endhead
Fixed single-tape Turing machine & Partial numerical functions; convergence
means accepting with a valid canonical output, while rejection, divergence,
and invalid acceptance are undefined. & Defines the Turing-computable class
under exactly this convergence convention. &
Definitions~\ref{def:p04-turing-machine} and
\ref{def:p04-partial-computable}; the next chapter supplies the index
realization. \\
\addlinespace
Partial recursive functions & Least partial-function class containing zero,
successor, and projections and closed under composition, primitive recursion,
and the definedness-sensitive \(\mu\)-operator. & Same partial numerical
functions as the fixed Turing-machine formalism. & Cited
Theorem~\ref{thm:p04-tm-partial-recursive};
\cite{Kleene1936,Kleene1952,Turing1937Lambda}. \\
\addlinespace
Shepherdson--Sturgis URM & Countably many \(\N\)-valued registers and finite
programs using exactly \(P,D,O,C,J,J\)-zero as specified in
Definition~\ref{def:p04-shepherdson-sturgis-urm}. & Computes exactly the
partial recursive numerical functions; equivalence with the preceding Turing
model follows through that shared class. & Primary URM construction and
two-symbol-machine reduction: \cite{ShepherdsonSturgis1963}; the final
TM--URM edge is transitive. \\
\addlinespace
Lambda definability & Fix Church numerals; on represented inputs, a term
represents a numerical function when its application to the input numerals
converts to the numeral for the value.  Comparison is by represented
functions, not raw-term equality. & Turing's historical theorem identifies the
corresponding computable, lambda-definable, and general recursive function
classes on positive integers. & \cite{Turing1937Lambda}; no modern
undefined-input convention is attributed to that paper. \\
\addlinespace
Effective calculability & Informal explicandum, not a fourth formal model or a
class fixed by a formal definition. & Church--Turing thesis pointer; no
equivalence edge is asserted from a theorem. &
\cite{Church1936,Turing1937Lambda,Soare1996ComputabilityRecursion}; see
Boundary~\ref{def:p04-church-turing}. \\
\bottomrule
\end{longtable}
\endgroup
\end{landscape}

\begin{proofboundary}{def:p04-church-turing}{Church--Turing thesis}
The \emph{Church--Turing thesis} identifies the informally understood
effectively calculable numerical functions with the Turing-computable
functions.  Church proposed general recursion and Turing analyzed machine
computation; the equivalence and robustness of the resulting formal models
support the identification
\cite{Church1936,Turing1937Lambda,Soare1996ComputabilityRecursion}.  They do
not prove it, because ``effective calculability'' is the informal explicandum,
not an independently formalized class quantified over by a theorem.  The
thesis stated here is not the claim that every physical process is Turing
computable, and formal-model equivalence does not imply equality of complexity
costs (Warning~\ref{def:p04-formalism-complexity}).
\index{Church--Turing thesis}
\end{proofboundary}

\section{Decidable and computably enumerable sets}
\label{sec:p04-decidable-ce}

\begin{hbdefinition}{def:p04-decidable-ce-set}{Decidable and computably enumerable sets}
For \(A\subseteq\N\), \(A\) is \emph{decidable} iff its characteristic
function \(\operatorname{char}_A:\N\to\{0,1\}\) is computable.  It is \emph{computably
enumerable}, abbreviated \emph{c.e.}, iff it is the domain of a partial
computable function.
\index{decidable set}
\index{computably enumerable set}
\end{hbdefinition}

\begin{hbdefinition}{def:p04-decidable-ce-relation}{Decidable and computably enumerable relations}
For \(R\subseteq\N^k\), \(R\) is \emph{decidable} iff its characteristic
function is computable.  It is \emph{computably enumerable} iff
\[
  \{\operatorname{code}_k(\vec n):\vec n\in R\}
\]
is a computably enumerable subset of \(\N\).  These terms classify relations;
for \(k=1\) they reproduce Definition~\ref{def:p04-decidable-ce-set}.  This
extension is used for computation-history and proof predicates.
\index{decidable relation}
\index{computably enumerable relation}
\end{hbdefinition}

\begin{hbdefinition}{def:p04-machine-decision-recognition}{Machine decision and recognition}
Let \(A\subseteq\N\).  A Turing machine \emph{decides} \(A\) iff it halts
on every unary input \(\operatorname{bin}(n)\), accepting exactly when
\(n\in A\) and rejecting exactly when \(n\notin A\).  It \emph{recognizes}
\(A\) iff it accepts exactly the inputs with \(n\in A\), with no halting
requirement outside \(A\).  Such machines are respectively a \emph{decider}
and a \emph{recognizer}.  These terms classify machine behavior, not sets or
functions; both are needed in the next two results, and they differ because a
recognizer may diverge outside \(A\).
\index{Turing machine!decider}
\index{Turing machine!recognizer}
\end{hbdefinition}

\begin{hblemma}{thm:p04-machine-decidability}{proved=annex:p04-machines.2}{Machine form of decidability}
A set \(A\subseteq\N\) is decidable iff some Turing machine decides \(A\).
\end{hblemma}

\begin{hbtheorem}{thm:p04-ce-equivalences}{sketched=annex:p04-machines.3}{Equivalent forms of computable enumerability}
For \(A\subseteq\N\), the following are equivalent:
\begin{enumerate}
  \item \(A\) is computably enumerable;
  \item some Turing machine recognizes \(A\);
  \item either \(A=\varnothing\), or \(A\) is the range of a computable
        function \(e:\N\to\N\); and
  \item there is a decidable relation \(R(x,k)\) such that
        \(x\in A\) iff \(\exists k\,R(x,k)\).
\end{enumerate}
The empty-set exception in the range clause is necessary because a total
function on \(\N\) has nonempty range.
\end{hbtheorem}

\begin{hbtheorem}{thm:p04-ce-complement}{proved=annex:p04-machines.4}{Complement criterion for decidability}
A set \(A\subseteq\N\) is decidable iff both \(A\) and
\(\N\setminus A\) are computably enumerable.
\end{hbtheorem}

% SOURCE ACCOUNTING -- RETIRED EXERCISE ex:turing-machines-01.
% The source exercise on closure of decidable sets under complement, finite
% union, and finite intersection, and of c.e. sets under finite union and
% finite intersection, is deliberately not printed: this handbook has no
% exercise layer.  Its warning not to assume c.e.-closure under complement is
% retained substantively by thm:p04-ce-complement.
